\section*{Chapter \ref{ch:pl:the-research-environment} --- The Research Environment}

\begin{solution}{exo:pl:the-research-environment:1}
The fourth code cell, with counter 3, is out of order (the cell above it has 5); counter 4 is missing (a re-run or a deleted cell). The
third cell ran after the fourth, and something else ran between them.
\end{solution}

\begin{solution}{exo:pl:the-research-environment:2}
No cell of the notebook binds \texttt{exclude}, above or below: it was bound by a cell that has been deleted. ``Used before defined'' is for a
name that a cell further down binds --- a rerun from the top fails there too, but the definition is still in the document.
\end{solution}

\begin{solution}{exo:pl:the-research-environment:3}
The five countries excluded by the deleted cell, and the years before 1950, which were still in the data when the result ran because the
cleaning cell ran afterwards. With the exclusion and all years the library gives 1.96; with every country from 1950, 2.44.
\end{solution}

\begin{solution}{exo:pl:the-research-environment:4}
Yes. A cell edited after it ran shows its new source with the old output; a cell that appends to a list, run once with consecutive counters and
then re-run after a restart that re-used counter numbers, can hide nothing in the counters; a name set at the prompt of an attached console
leaves no trace. Only the rerun is conclusive.
\end{solution}

\begin{solution}{exo:pl:the-research-environment:5}
Both answer a question; they are different questions --- the typical country's experience against the typical country-year's --- and the
difference (2.44 against 2.90) is larger than many effects reported. A function that hides the choice publishes one answer as if it were the
only one; as an argument, it is stated in the call and in the report. Report the one that matches the question, and the other beside it.
\end{solution}

\begin{solution}{exo:pl:the-research-environment:6}
A lock is enough when the result depends only on the interpreter and the Python packages it pins, on a known operating system; a container is
needed when it depends on system libraries, compilers or tools outside the lock (\omterm{def:pl:native-extensions:extension}{native extensions}, database clients), or must run on machines
you do not control.
\end{solution}

\begin{solution}{exo:pl:the-research-environment:7}
On the fixed notebook the check has nothing to compare: the fix cleared the saved outputs, as it should, since they came from a history that no
longer exists. On the planted notebook the rerun stops at the fourth code cell, so the result cell's saved 1.96 is never compared with anything.
It is harder than comparing the last output because the rerun's output must be attributed to cells (markers printed between them), some outputs
legitimately differ (times, memory addresses, the last digits of floating-point sums), and rich outputs such as plots need their own comparison.
\end{solution}

\begin{solution}{exo:pl:the-research-environment:8}
Version control stores the document, with its saved outputs; it does not store the kernel's state that produced them, and a reviewer reads the
document. Reproducible means rerun: from a clean checkout, in the locked environment, on the snapshot of the data, from the top, with the same
result. Neither the commit nor the approval did that.
\end{solution}

\begin{solution}{pb:pl:the-research-environment:1}
\begin{enumerate}
\item Its source, its saved outputs, and the kernel's execution counter when it last ran (an integer, or null).
\item Run: load, select above 90\%, the deleted exclusion, the country selection, the result, then the cleaning; shown: load, cleaning,
selection, country selection, result.
\item 1.96.
\item A cell run out of order, a cell re-run, a cell edited or deleted after it ran (and names set at the prompt).
\item Of 863\,878 attempted executions of valid notebooks, 24.11\% ran without errors and 4.03\% reproduced their results.
\item Out of order (cell 3 ran before cell 2), a skipped counter (3), and a name defined nowhere (\texttt{exclude}); from the JSON alone.
\item It fails: \texttt{NameError} on \texttt{exclude}.
\item The skipped counter and the undefined name are the deleted exclusion; the out-of-order counter says the displayed data are not those the
result used.
\item Edits after a run, re-runs that leave consecutive counters, and anything done outside the notebook.
\item So that nothing of the author's session --- names, imports, working directory --- leaks into the check.
\item The threshold, the first year, the excluded countries and the weighting.
\item 1.96 (five countries excluded, all years), 2.44 (all countries from 1950, by country), 2.90 (all countries from 1950, by country-year).
\item 204 of 455.
\item The checklist: no \omterm{def:pl:the-research-environment:hidden}{hidden state}, rerun succeeds and reproduces, tests, environment against the lock, data against the snapshot.
\item The library and every caller change together and are tested together.
\item \textbf{Named result.} The notebook shows 1.96; rerun from the top it fails on a name defined only by a deleted cell; the audit finds that
cell (a skipped counter and a name defined nowhere) and a cleaning cell run after the result; corrected, the calculation gives 2.44, which the
promoted library reproduces from a clean directory with the data snapshot and environment lock checked.
\item The promoted calculation, a rerun from a clean checkout that prints it, the lock and the data snapshot, and the choices stated in the
call.
\item When the result depends on anything outside the lock, or must run on machines I do not control.
\item The question, the choices (threshold, years, countries, weighting), the command that recomputes the result, the code and data hashes, and
the alternatives tried.
\item That anyone can recompute it from the top, from a clean checkout, and get the same number.
\end{enumerate}
\end{solution}

\begin{solution}{iq:pl:the-research-environment:1}
Rerun it from the top in a fresh kernel; read its execution counters for skips and out-of-order runs; look for names defined nowhere; check the
data version and the package versions against what the author used.
\iqlookfor{Rerun first, then \omterm{def:pl:the-research-environment:hidden}{hidden state}, data and environment.}
\end{solution}

\begin{solution}{iq:pl:the-research-environment:2}
As soon as it is used twice, feeds anything others rely on, or its choices matter: then it needs arguments for its choices, tests, and review.
\iqlookfor{Reuse and consequence, not size.}
\end{solution}

\begin{solution}{iq:pl:the-research-environment:3}
The whole environment below the Python packages --- operating system, system libraries, compilers --- frozen and portable. It costs build time,
image storage and maintenance, and an image can still hide what is in it unless it is built from a lock.
\iqlookfor{Freezing versus describing; build from a lock.}
\end{solution}

\begin{solution}{iq:pl:the-research-environment:4}
Produce every number by a command in the repository that runs from a clean checkout on a data snapshot in a locked environment, and record in the
report the commit, the snapshot hash and the command; test in continuous integration that the commands still reproduce the numbers.
\iqlookfor{Numbers as outputs of versioned commands.}
\end{solution}

\begin{solution}{iq:pl:the-research-environment:5}
Notebooks for exploration, a shared tested \omterm{def:pl:the-research-environment:library}{research library} in one repository with review, environments from a lock (containers where needed),
data read through the platform's snapshots, an audit and rerun gate before any result is shared, and a research log with the commands that
reproduce each result.
\iqlookfor{Exploration free, results gated.}
\end{solution}
